\section{Future Directions}
\label{sec:future}

This section specifies deferred capabilities at architectural depth --- enough
to confirm that the v1.0 design does not foreclose them, and to identify the
specific v1.0 decisions that were made to keep them open.

\subsection{Video and Real-Time Emulation}

\paragraph{What changes} The pipeline becomes a per-frame operation at 24, 25,
or 30\,fps, with a hard budget of 33\,ms shared with capture, encode, and audio.
The pointwise chain is already a single LUT fetch and is essentially free. The
binding costs are halation (0.44\,ms at preview resolution on A17, but video is
captured at higher resolution) and grain.

\paragraph{The hard problem} Temporal grain coherence. Grain must be
\emph{uncorrelated} between frames --- real grain is a new realization each
frame --- but the underlying image is correlated, so naive per-frame
regeneration produces the correct result. The genuine difficulty is that video
codecs allocate bits to changing regions, and uncorrelated full-frame noise
consumes the entire bitrate. Professional practice is to grade with grain and
deliver at a bitrate sufficient to carry it; on a phone this is a real
constraint.

The specified approach is to record the grain \emph{parameters} in the container
alongside a clean encode, and synthesize grain at playback within the
application, with a baked-in option for export. This is architecturally
identical to the film grain synthesis facility standardized in AV1, and the v1.0
decision that makes it possible is that grain is fully determined by
(seed, parameters, film coordinate) --- Section~\ref{sec:grain} --- so it can be
regenerated exactly from a handful of bytes.

\paragraph{v1.0 decisions that preserve this} Resolution-independent,
seed-deterministic grain; physical-unit specification of all spatial phenomena;
the render graph's resolution-relative descriptors.

\subsection{Log Input and Professional Interchange}

Apple Log and Blackmagic Log inputs require only an additional input transform:
a log-to-linear decode followed by the existing $\mathbf{M}_{\mathrm{in}}$. The
pipeline is already scene-referred, which is the property that makes this a
small change rather than a redesign.

Export of the baked pointwise chain as a \texttt{.cube} 3D LUT is likewise
nearly free --- the LUT already exists (Section~\ref{sec:lutbake}) and needs
only serialization plus a documented shaper. This is a high-value, low-cost
feature for the P2 persona: it lets a colorist take an \EM{} grade into
Resolve directly. It is scheduled for 1.2.

The honest caveat, which must be surfaced in the UI when exporting a LUT: the
spatial and stochastic stages cannot be baked. A \texttt{.cube} export carries
the color science and none of the halation, grain, or interlayer effects, and
telling the user so is more valuable than a silently incomplete export.

\subsection{ACES Integration}

The working space is already ACEScg. Full ACES integration adds: IDT selection
for non-Apple sources, an ACES container output (OpenEXR with the appropriate
metadata), and --- the substantive part --- a decision about where the film
model sits relative to the ACES viewing transform.

The correct answer is that \EM{}'s negative-plus-print chain \emph{is} an
output transform: it takes scene-referred data to display-referred, which is
exactly the RRT+ODT's role. \EM{} therefore replaces the ACES output transform
rather than composing with it, and an ACES-integrated build must expose that
choice explicitly rather than silently double-transforming. This is stated here
because it is the mistake that ACES integrations most commonly make.

\subsection{Node-Based Pipeline Editing}

The render graph is already a DAG with declared reads, writes, and parameter
dependencies (Section~\ref{sec:rendergraph}). Node-based editing requires: a
serialized graph topology in the recipe rather than a fixed chain; a UI for
constructing it; validation that user-constructed graphs are acyclic and
domain-consistent (a pass consuming density must not be fed linear exposure);
and a defined behavior for graphs whose stages appear in a physically impossible
order.

That last item is the interesting design problem, and it is a direct conflict
with the parameter-honesty principle. The proposed resolution is that arbitrary
ordering is permitted but non-physical orderings are marked in the interface,
with an explanation of what the physical ordering would be. Freedom with
labelling, rather than prohibition or silence.

\subsection{User-Authored and Community Profiles}

The stock profile format is already a versioned JSON parameter block
(Section~\ref{sec:schema}). Opening it requires: a profile editor exposing the
parameters that are currently profile-only ($\mathbf{A}$, $\mathbf{C}$, grain
$\sigma_1$/$\chi$, halation $\ell$ and $\boldsymbol{\beta}$); a validator
enforcing the well-formedness condition \eqref{eq:wellformed} and the
monotonicity requirement AC-6; and a distribution mechanism.

The genuinely valuable version of this feature is not authoring by eye but
\emph{fitting from measurement}: a user photographs a step wedge on their own
film, scans it, and the application fits $\theta$. That is a substantially more
defensible feature and it aligns with the product's identity. It requires a
guided capture flow and a robust fitter, and it is the correct 3.0 headline
feature.

\subsection{Machine Learning: Where It Helps and Where It Does Not}

The original proposal contemplates machine learning for exposure recommendation
and for estimating spectral film responses from scanned negatives. These are
very different propositions and should be evaluated separately.

\paragraph{Estimating model parameters from scans --- promising} Fitting
$\theta$, $\mathbf{A}$, $\mathbf{C}$, and the grain parameters from a set of
scanned reference images is a well-posed inverse problem. Because the forward
model of Sections~\ref{sec:negative}--\ref{sec:grain} is differentiable by
construction (this is why the softplus formulation was chosen over a spline),
gradient-based optimization applies directly, and the "learning" is parameter
estimation within a physical model rather than a black box. The result remains
interpretable, editable, and validatable against AC-1. This is the right use.

\paragraph{End-to-end learned film emulation --- rejected} A network mapping
RAW to film-look output would discard every property the product is built on:
interpretability, parameter honesty, controllability, and the ability to
validate. It would also be immediately commoditized. This direction is
explicitly not pursued.

\paragraph{Exposure and stock recommendation --- modest} A small on-device
model suggesting exposure compensation and stock from scene content is a
convenience feature with real but bounded value. It should be additive and
always overridable, and it should present its reasoning in the product's own
vocabulary ("this scene has a 9-stop range; a stock with a longer straight line
will hold both ends").

\subsection{Full Spectral Rendering}

The reduced spectral strategy of Section~\ref{sec:colorspace} bakes crosstalk at
profile time against a reference illuminant. A full spectral path would carry
$N \approx 31$ wavelength samples through the exposure and density stages,
integrating against measured layer sensitivities and dye transmittances.

The cost is roughly $10\times$ in the density stage --- which, per
Table~\ref{tab:previewperf}, is a small fraction of the total, so the cost is
more tolerable than it first appears. The real obstacle is data: measured
spectral sensitivities and dye densities are not published for most stocks and
would require spectrophotometric measurement of physical samples.

The value would be correct rendering under unusual illuminants and correct
prediction of crosstalk from first principles rather than by fitted matrix.
This is a research direction rather than a product feature, appropriate for a
later phase once the profile measurement apparatus already exists for other
reasons.

\subsection{Cross-Device Continuity}

The macOS companion shares \texttt{EmulsionCore} and \texttt{EmulsionRender}
unchanged; only \texttt{EmulsionCapture} and the UI layer differ. This is the
direct payoff of the module layering in Table~\ref{tab:layers} and the reason
that layering was specified as strictly as it was.

Synchronization carries recipes and library metadata, not images. A recipe is
under a kilobyte; syncing recipes is trivial and syncing 48\,MP originals is a
different product with different economics. The sync model is therefore:
recipes and metadata via CloudKit, images by user action only. This keeps the
privacy commitment of CR-02 substantially intact --- the default remains that
no image leaves the device.
